\itemhead{习题4.1 · 第5题（教材p.104）}{XJ-B1-C04-HW-P104-Y1201}
\label{XJ-B1-C04-HW-P104-Y1201}
{\footnotesize 湘教版教材p.104 / PDF第111页 / 习题4.1 / 第5题}
\field{原题}{计算：
\[\sqrt[3]{2-\frac{62}{27}}+\sqrt{\left(-3\frac23\right)^2}-3\div16^{-0.75}+\sqrt[5]2\cdot\left(4^{-1/5}\right)^{-2}.\]}
\field{参考答案}{$-19$。四部分依次为 $-2/3,11/3,-24,2$。}
{\small \field{关联知识}{P02 有理数、实数指数幂与运算（教材帮p.144、145 / PDF153、154）}
\field{方法与易错点}{统一为指数幂，分开处理系数与各底数的指数；最终复查定义条件及数值精度。 易错：负指数误作负数；幂的幂指数相加；约分丢禁点；过早四舍五入。}
\field{答案依据}{答案册相关题组：PDF第31页（册内页30）；按当前题库原题逐项重算，不按旧版题号机械对套。}}
